Graphs are the natural data structure of relations. Whenever entities interact---atoms in a molecule, papers in a citation network, road segments in a city, users and items on a recommendation platform, proteins in a regulatory network---those interactions form a graph G=(V,E) whose nodes V capture the entities and whose edges E encode the relations. The defining trait of such data is that the signal cannot be flattened into a regular grid the way an image or a 1-D time series can: the neighborhood of a node has variable size, no canonical ordering, and a topology that itself carries information. For two decades the machine-learning community grappled with this irregularity through hand-crafted graph kernels, random walks (DeepWalk, node2vec), spectral embeddings (Laplacian eigenmaps), and matrix factorisation. Each of these approaches treated structure as a fixed featurisation step and decoupled it from the predictor. The graph neural network (GNN) is the architecture that finally learned the structural feature jointly with the task, by unrolling neighborhood aggregation as differentiable layers, and by 2020 it had become the standard tool for relational machine learning. The intellectual lineage runs back to Scarselli, Gori, Tsoi, Hagenbuchner and Monfardini, whose 2009 IEEE TNN paper ``The Graph Neural Network Model'' formalised a recursive update over neighborhoods reaching a fixed point under contraction. Their model was theoretically appealing but practically constrained b\ldots{}
